% SGA 3, Expose III, Loop 1 English reconstruction.
% Scope: Lemma 0.1.11 through Proposition 0.2, stopping before Remark 0.3.
% Authority: Polo--Gille Exp3-14oct24.pdf, component-local pages 12--14
% (combined-reader pages 122--124; printed pages 112--114, with original
% pagination 88--90 also shown by the re-edition). Force-OCR is a drafting
% locator; English wording was compared with Jacob Reinhold's CC BY 4.0
% Markdown lineage. The PDF controls formulas, notes, and diagram images.

\setcounter{footnote}{17}

% Source page: Exp3 component-local p. 12; combined-reader p. 122.
\medskip
\noindent\SGAInlineRefTarget{sga3:iii:lemma:0-1-11}\textbf{引理 0.1.11.}
\textit{设 \(X,X'\) 是两个 \(S\)-概形. 有一个交换图}

\SGAThreeNativeDiagram[0.58\textwidth]
  {figures/exp3/diagram_015_exp3_p012_lemma0111_product_square.tex}
  {\(L_X\) 与 \(X^+\) 关于乘积的相容性.}
  {fig:exp3-lemma-0111-product}
  {Polo--Gille Exp3-14oct24.pdf, 组件第 12 页; 合卷第 122 页}

\textit{证明.}\footnote{编者注: 证明作了展开.}
对每个 \(S\)-概形 \(Y\),
\[
\begin{aligned}
 \operatorname{Hom}_S(Y,X^+\times_S X'^+)
 &=
 \operatorname{Hom}_S(Y,X^+)
 \times\operatorname{Hom}_S(Y,X'^+)\\
 &\simeq
 \operatorname{Hom}_S(Y_{\mathscr J},X)
 \times\operatorname{Hom}_S(Y_{\mathscr J},X')\\
 &=
 \operatorname{Hom}_S\bigl(Y,(X\times_S X')^+\bigr).
\end{aligned}
\]
因此
\[
 X^+\times_S X'^+\simeq(X\times_S X')^+.
\]

现设 \(Y\) 是 \(X^+\times_S X'^+\) 上的概形, 由态射
\[
 h:Y_{\mathscr J}\longrightarrow X\times_S X'.
\]
表示. 设 \(p,q\) 为从 \(X\times_S X'\) 到 \(X,X'\) 的投影,
并令 \(f=p\circ h\), \(g=q\circ h\). 由于
\[
 \Omega^1_{(X_0\times_{S_0}X'_0)/S_0}
 \simeq
 p_0^*(\Omega^1_{X_0/S_0})
 \oplus q_0^*(\Omega^1_{X'_0/S_0})
\]
(EGA~IV\textsubscript{4}, 16.4.23), 有一个自然同构
\[
\begin{aligned}
 &\operatorname{Hom}_{\mathscr O_{Y_0}}
   \bigl(f_0^*(\Omega^1_{X_0/S_0}),\mathscr J\mathscr O_Y\bigr)
 \times
 \operatorname{Hom}_{\mathscr O_{Y_0}}
   \bigl(g_0^*(\Omega^1_{X'_0/S_0}),\mathscr J\mathscr O_Y\bigr)\\
 &\qquad\simeq
 \operatorname{Hom}_{\mathscr O_{Y_0}}
 \bigl(h_0^*(\Omega^1_{(X_0\times_{S_0}X'_0)/S_0}),
       \mathscr J\mathscr O_Y\bigr).
\end{aligned}
\]
换言之,
\[
 L_X(Y)\times L_{X'}(Y)\simeq L_{X\times_S X'}(Y).
\]

\medskip
\noindent\SGAInlineRefTarget{sga3:iii:remark:0-1-12}\textbf{注 0.1.12.}
\footnote{编者注: 此注及其后的推论为新增内容.}
设 \(\mathscr C\) 是一个对纤维积稳定的范畴, \(S\) 是
\(\mathscr C\) 的对象, \(T_1,T_2\) 是 \(S\) 上的两个对象.
对 \(i=1,2\), 设 \(L_i\) 和 \(X_i\) 是 \(T_i\) 上的两个对象,
如下图所示.

\SGAThreeNativeDiagram[0.52\textwidth]
  {figures/exp3/diagram_016_exp3_p012_remark0112_fiber_objects.tex}
  {\(T_i\) 上的对象 \(L_i,X_i\), 以及 \(S\) 上的 \(T_i\).}
  {fig:exp3-remark-0112-objects}
  {Polo--Gille Exp3-14oct24.pdf, 组件第 12 页; 合卷第 122 页}

于是有一个自然同构
\[
 (L_1\times_{T_1}X_1)\times_S(L_2\times_{T_2}X_2)
 \simeq
 (L_1\times_S L_2)
 \times_{T_1\times_S T_2}(X_1\times_S X_2).
\]
因此, 由前一引理得到:

% Source page: Exp3 component-local p. 13; combined-reader p. 123.
\medskip
\noindent\SGAInlineRefTarget{sga3:iii:corollary:0-1-13}\textbf{推论 0.1.13.}
\textit{设 \(X_1,X_2\) 是两个 \(S\)-概形. 有一个由同构构成的交换图:}

\SGAThreeNativeDiagram[0.96\textwidth]
  {figures/exp3/diagram_017_exp3_p013_corollary0113_isomorphisms.tex}
  {\SGARefLink{sga3:iii:corollary:0-1-13}{推论 0.1.13} 中的乘积同构.}
  {fig:exp3-corollary-0113}
  {Polo--Gille Exp3-14oct24.pdf, 组件第 13 页; 合卷第 123 页}

现在可以陈述:

\medskip
\noindent\SGAInlineRefTarget{sga3:iii:proposition:0-2}\textbf{命题 0.2.}
\textit{对每个 \(S\)-概形 \(X\), 都可定义 \(X^+\)-交换群
\(L_X\) 在 \(X^+\)-对象 \(X\) 上的一个左作用,
具有以下性质:}
\begin{enumerate}
\item[\SGAInlineRefTarget{sga3:iii:proposition:0-2:item:i}(i)]
\textit{此作用使 \(X\) 在 \(X^+\) 上关于 \(L_X\) 为形式主齐性的;
即态射}
\[
 L_X\times_{X^+}X\longrightarrow X\times_{X^+}X
\]
\textit{是 \(X^+\)-函子同构;}

\item[\SGAInlineRefTarget{sga3:iii:proposition:0-2:item:ii}(ii)]
\textit{此作用关于 \(S\)-概形 \(X\) 具有函子性:
对每个 \(S\)-态射 \(f:X\to W\), 下图交换;}

\SGAThreeNativeDiagram[0.55\textwidth]
  {figures/exp3/diagram_018_exp3_p013_proposition02_functoriality_square.tex}
  {\SGARefLink{sga3:iii:proposition:0-2:item:ii}{命题 0.2(ii)} 中作用的函子性.}
  {fig:exp3-proposition-02-functoriality}
  {Polo--Gille Exp3-14oct24.pdf, 组件第 13 页; 合卷第 123 页}

\item[\SGAInlineRefTarget{sga3:iii:proposition:0-2:item:iii}(iii)]
\textit{此作用与纤维积可交换: 对任意 \(S\)-概形 \(X_1,X_2\),
下图交换.}
\end{enumerate}

\SGAThreeNativeDiagram[0.99\textwidth]
  {figures/exp3/diagram_019_exp3_p013_proposition02_product_diagram.tex}
  {作用与纤维积的相容性.}
  {fig:exp3-proposition-02-products}
  {Polo--Gille Exp3-14oct24.pdf, 组件第 13 页; 合卷第 123 页}

\textit{证明.}\footnote{编者注: 为计入前面新增的内容,
对原文作了修改和展开; 特别地, 这些新增内容详细纳入了原版第 89 页的内容.}
\SGARefLink{sga3:iii:proposition:0-2:item:i}{(i)} 和
\SGARefLink{sga3:iii:proposition:0-2:item:ii}{(ii)} 分别由
\SGARefLink{sga3:iii:corollary:0-1-9}{推论 0.1.9} 和
\SGARefLink{sga3:iii:corollary:0-1-10}{0.1.10} 得到.
为证明 \SGARefLink{sga3:iii:proposition:0-2:item:iii}{(iii)},
对每个 \(S\)-概形 \(X\), 记
\[
 P(X)=L_X\times_{X^+}X
\]
将 \SGARefLink{sga3:iii:proposition:0-2:item:ii}{(ii)} 应用于投影
\[
 p_i:X_1\times_S X_2\longrightarrow X_i
\]
得到以下交换方形.

\SGAThreeNativeDiagram[0.55\textwidth]
  {figures/exp3/diagram_020_exp3_p013_proposition02_projection_square.tex}
  {每个投影 \(p_i\) 诱导的方形.}
  {fig:exp3-proposition-02-projection-square}
  {Polo--Gille Exp3-14oct24.pdf, 组件第 13 页; 合卷第 123 页}

% Source page: Exp3 component-local p. 14; combined-reader p. 124.
对 \(i=1,2\), 由这些方形得到交换方形

\SGAThreeNativeDiagram[0.55\textwidth]
  {figures/exp3/diagram_021_exp3_p014_proposition02_final_square.tex}
  {\SGARefLink{sga3:iii:proposition:0-2:item:iii}{命题 0.2(iii)} 证明中的乘积方形.}
  {fig:exp3-proposition-02-final-square}
  {Polo--Gille Exp3-14oct24.pdf, 组件第 14 页; 合卷第 124 页}

将此方形与 \SGARefLink{sga3:iii:corollary:0-1-13}{推论 0.1.13}
结合, 可知其竖直箭头是同构, 并恰好得到
\SGARefLink{sga3:iii:proposition:0-2:item:iii}{(iii)} 中所述的交换图.

% Continue with Remark 0.3.
% Translation statistics: source_lines=179; output_lines=182; inline_math=48; display_math=10; equation_environments=0; footnotes=3.
% Review: complete assigned component; pending_TODO=0; math-token multiset, equation blocks, targets, references and environments checked.
